\section{Simultaneous precision and recall targets}
An adaptive follow-up specifies precision $>82\%$, overall recall $>90\%$,
and distant-turn recall $>80\%$ simultaneously. These are point-estimate targets,
not confidence-bound requirements. We retain the 104/52 development subdivision
and seal confirmation. Optuna selection uses only calibration; one frozen
candidate is scored per sweep. These reused-development, agent-labeled outcomes
cannot establish independent generalization.

The first 256-trial TPE sweep adjusts Code 4 2048-dimensional BIRCH radius,
membership cutoff, relative-best margin and view. A 512-trial follow-up adds
convex representation ensembles: concatenate $\sqrt{w}$ times normalized Code
vectors and $\sqrt{1-w}$ times normalized Voyage vectors. The resulting cosine
is the weighted sum of the model cosines. Five weights are registered.
A final 768-trial sweep substitutes average-link, complete-link or online leader
partitioning, retaining overlapping centroid memberships and distinct-turn
support. We optimize the minimum ratio to the three targets until feasible;
among feasible calibration settings, precision is maximized.

\begin{center}\small
\begin{tabular}{lrrrrrr}
\toprule
& \multicolumn{3}{c}{Calibration} & \multicolumn{3}{c}{Development evaluation}\\
Sweep & P & R & Distant R & P & R & Distant R\\
\midrule
Code BIRCH &81.47&92.42&82.14&77.83&89.12&77.78\\
BIRCH ensemble &83.52&90.46&86.90&76.60&93.26&88.89\\
Family comparison &83.11&90.22&88.10&76.72&92.23&88.89\\
Group refinement &83.52&90.46&86.90&76.60&93.26&88.89\\
\bottomrule
\end{tabular}
\end{center}
The latter two calibration candidates pass all targets, but evaluation precision
fails. The BIRCH ensemble selects $w=.25$, radius .45, membership cosine .6375,
relative margin .20 and chunk view. The family sweep selects complete-link at
partition cosine .60 with the same weight and membership settings. Each recovers
48 of 54 distant related evaluation pairs, versus 38 for the frozen reference,
but admits 55 and 54 false positives respectively, versus 33. Their evaluation
F1 scores are 84.11 and 83.76, with exploratory paired gain intervals
$[-3.86,+1.11]$ and $[-5.27,+1.78]$ percentage points against the reference.
Neither demonstrates a balanced-quality improvement.

Calibration contains 409 related and 188 unrelated scored pairs; evaluation
contains 193 and 103 respectively. We exclude 68/50 ambiguous judgments.
Distant denominators are only 84/54. This finite evidence should not be treated
as a corpus-wide performance guarantee. The precision generalization gap
motivates a distinct weak-link check rather than more radius-only tuning.
An additional 33-trial registered refinement splits overlapping centroid groups
using complete-link on unsmoothed joint-turn vectors. Low cosine floors do not
change calibration outcomes; stronger floors lose recall. Selection retains
disabled refinement, exactly replaying the BIRCH ensemble. This row is a control
replay rather than a distinct quality improvement.
All 1,569 trials use cached vectors, cost no additional inference, and preserve
thread-local processing. Driver times are 20.6, 52.1 and 62.1 seconds, excluding
final SQLite persistence; refinement takes 21.9 seconds. Optuna trials and MLflow
aggregate records remain local. The drivers pass Ruff, ty and all 168 Python tests, with
regressions covering strict target selection, convex cosine geometry, distinct
turn support and overlapping refinement. Production and the
frozen baseline remain unchanged.
